% ============ 第六章 Abel 数域的类数公式 (原书 pp.206-234) ============
\chapter{Abel 数域的类数公式}\label{ch:6}

我们现在展示解析理论在研究代数数域类数问题中的应用. 解析理论与类数问题的联系已经在第四章中给出, 那就是:
\[
  \lim_{s\to1}(s-1)\zeta_K(s)=\frac{2^{r_1}(2\pi)^{r_2}R_Kh_K}{w_K|d(K)|^{1/2}}.
\]
此式右边出现的数域 $K$ 中的诸量 $r_1,r_2,w_K,|d(K)|^{1/2}$ 通常容易求得. $R_K$ 与基本单位组有关, 是难求的量. 对于此式的左边, 我们则希望 $\zeta_K(s)$ 有更好的表达式. 高斯对于二次域作了这件事, 从而得到了二次域类数相当简单的公式. 后来, 在研究费马问题的推动之下, Kummer 对于分圆域也给出了类数解析公式. 到了 $20$ 世纪 $30$ 年代之后, Hasse 对于一般的 Abel 数域给出类数解析公式, 并且由此相当精细地研究了 Abel 域类数的各种问题. 其研究成果集中总结在他于 $1952$ 年所写的 ``关于 Abel 域的类数'' 一书中.

% ============ 第六章 §1 Hasse 类数公式 (原书 pp.206-216) ============
\section{Hasse 类数公式}\label{sec:ch6-s1}

对于 Abel 数域 $K$, $\zeta_K(s)$ 应当有好的表达式, 这一点是不应当感到很奇怪的. 因为根据 Hronecker-Weber 定理, 每个 Abel 数域均是分圆域的子域, 而分圆域中有很简单的素理想分解规则. 设 $\Q(\zeta_m)(m\not\equiv 2\ (\mathrm{mod}\ 4))$ 是包含 Abel 数域 $K$ 的最小分圆域, 即 $m=\operatorname{cond}(K)$（域 $K$ 的导子）. 我们有如下的伽罗瓦对应:
\[
  \begin{array}{cc}
    \Q(\zeta_m) & G=\Gal(\Q(\zeta_m)/\Q)\\
    \mid & \mid\\
    K & H\\
    \mid & \mid\\
    \Q & \{1\}
  \end{array}
\]
由于
\[
  G\cong(\Z/m\Z)^\times,\qquad\sigma_a\mapsto a\ (\mathrm{mod}\ m),\quad (a,m)=1,\quad \sigma_a(\zeta_m)=\zeta_m^a.
\]
从而 $H$ 同构于 $(\Z/m\Z)^\times$ 的一个子群（我们今后常常把 $H$ 和这个子群等同起来）, 而 $\Gal(K/\Q)$ 同构于商群 $(\Z/m\Z)^\times/H$（我们也常把这两者等同起来）. 有限 Abel 群 $\Gal(K/\Q)=(\Z/m\Z)^\times/H$ 的每个特征 $\chi$ 也叫作是域 $K$ 的特征, 而 $(\Z/m\Z)^\times/H$ 的特征群也叫作是 $K$ 的特征群, 表示成 $\widehat{K}$. 于是 $\widehat{K}$ 实际上是由全部在 $H$ 上平凡的模 $m$ 的 D-特征所构成的.

设 $\chi$ 为模 $m$ 的 D-特征, 则 $\chi$ 是由惟一决定的一个模 $m'$ 的本原 D-特征所诱导出来的, 其中 $m'=\operatorname{cond}(\chi)$, $m'\mid m$. 我们将诱导出 $\chi$ 的这个本原 D-特征表示成 $\chi^*$.

如果 $\Q\subseteq K\subseteq L\subseteq\Q(\zeta_m)$, 而 $E_K$ 和 $E_L$ 分别为 $G=\Gal(\Q(\zeta_m)/\Q)$ 对于 $K$ 和 $L$ 的固定子群, 则 $E_K\subseteq E_L$, 即伽罗瓦对应 $K\leftrightarrow E_K$ 是反序的（序指的是集合之间的包含关系）. 由于 $\Gal(K/\Q)\simeq G/E_K$, $\Gal(L/\Q)\cong G/E_L$, 可知它们的特征群之间有关系 $\widehat{K}\subseteq\widehat{L}$. 所以对应 $K\mapsto\widehat{K}$ 是保序的.

群 $G$ 同构于 $(\Z/m\Z)^\times$. 若 $m=p_1^{e_1}\cdots p_r^{e_r}$ 是 $m$ 的素因子分解式, 则有同构
\[
  (\Z/m\Z)^\times\cong(\Z/p_1^{e_1}\Z)^\times\times\cdots\times(\Z/p_r^{e_r}\Z)^\times\qquad(\text{直积}),
\]
所以每个模 $m$ 的 D-特征惟一分解为 $\chi=\chi_{p_1}\cdots\chi_{p_r}$, 其中 $\chi_{p_i}$ 是模 $p_i^{e_i}$ 的特征（$1\leqslant i\leqslant r$）. 易知 $f_\chi=f_{\chi_{p_1}}\cdots f_{\chi_{p_r}}$, 这里 $f_\chi$ 表示特征 $\chi$ 的导子. 如果 $p$ 不为 $p_i(1\leqslant i\leqslant r)$, 则 $\chi_p$ 是主特征. 现在对 Abel 数域 $K$ 和每个素数 $p$, 定义如下一些特征集合:
\[
  \widehat{K}_p=\{\chi_p:\chi\in\widehat{K}\},\qquad
  Y=Y(K,p)=\{\chi\in\widehat{K}:\chi^*(p)\neq 0\},\qquad
  Z=\{\chi\in Y:\chi^*(p)=1\}.
\]
易知它们均形成特征群, 并且 $Z\subseteq Y\subseteq\widehat{K}$. 下面引理表明: 这些特征群可以刻画 $p$ 在 $K$ 中的分解型式.

\begin{lemma}\label{lem:ch6-1}
以 $e,f,g$ 分别表示 $p$ 在 Abel 数域 $K(\subseteq\Q(\zeta_m))$ 中的分歧指数, 剩余类域次数和分解次数, 则
\begin{enumerate}[(1)]
  \item $e=|\widehat{K}_p|$, 并且 $p$ 在 $K$ 中分歧当且仅当存在 $\chi\in\widehat{K}$, 使得 $\chi^*(p)=0$（即 $p\mid f_\chi$）.
  \item $Y$ 和 $Z$ 分别是 $p$ 对 $K$ 的惯性域和分解域的特征群. 于是
    \[
      e=[\widehat{K}:Y],\qquad f=[Y:Z],\qquad g=|Z|.
    \]
\end{enumerate}
\end{lemma}

\begin{proof}
(1) 设 $m=p^ln$, $(p,n)=1$, $l\geqslant 0$, 则有如下的域扩张图表:
\[
  \begin{array}{cc}
    \Q(\zeta_m)=\Q(\zeta_{p^l})\Q(\zeta_n) & (\Z/n\Z)^\times\times\widehat{K}_p\\
    \mid & \mid\\
    L=\Q(\zeta_n)K & \widehat{L}\\
    \mid & \mid\\
    K & \widehat{K}\\
    \mid & \mid\\
    \Q & \{1\}
  \end{array}
\]
图中每个域的特征群均看成是 $\Q(\zeta_m)$ 的特征群（即模 $m$ 的 D-特征群）的子群. 由于 $L$ 是 $\Q(\zeta_n)$ 和 $K$ 的合成, 而域到此域特征群之间的一一对应是保序的, 所以 $\widehat{L}$ 是由 $\Q(\zeta_n)$ 和 $\widehat{K}$ 生成的特征群. 于是
\[
  \widehat{L}=\langle(\Z/n\Z)^\times,\widehat{K}\rangle=(\Z/n\Z)^\times\ \text{和}\ \widehat{K}_p\ \text{的直积}.
\]
以 $F$ 表示对应于 $\widehat{K}_p$ 的域, 即 $\widehat{F}=\widehat{K}_p$, 则 $L=\Q(\zeta_n)F$, 并且 $F\subseteq\Q(\zeta_{p^l})$. 对于 $\Q(\zeta_m)$ 的每个子域 $M$, 以 $e(M)$ 表示 $p$ 在 $M$ 中的分歧指数, 则
\begin{align*}
  e=e(K)&=e(L)\qquad(\text{因为 } e(\Q(\zeta_n))=1\ \text{而 } L=\Q(\zeta_n)K)\\
  &=e(F)\qquad(\text{因为 } L=\Q(\zeta_n)F\ \text{而 } e(\Q(\zeta_n))=1)\\
  &=[F:\Q]\qquad(\text{因为 } F\subseteq\Q(\zeta_{p^l}),\ \text{从而 } p\ \text{在 } F\ \text{中完全分歧})\\
  &=|\widehat{F}|=|\widehat{K}_p|.
\end{align*}
进而, $p$ 在 $K$ 中分歧 $\Leftrightarrow|\widehat{K}_p|=e\geqslant 2\Leftrightarrow$ 存在 $\chi\in\widehat{K}$, 使得 $\chi_p\neq\chi_0\Leftrightarrow$ 存在 $\chi\in\widehat{K}$, 使得 $\chi^*(p)=0$. 这就证明了 (1).

(2) 设 $M$ 为 $K$ 的任一子域. 由 (1) 可知
\[
  p\ \text{在 } M\ \text{中不分歧}\Leftrightarrow\ \text{对每个 } \chi\in\widehat{M},\ \text{均有 } \chi^*(p)\neq 0
  \Leftrightarrow\widehat{M}\subseteq Y.
\]
这就表明 $Y$ 对应于 $p$ 在 $K$ 中的最大不分歧子域, 即对应于 $p$ 在 $K$ 中的惯性域 $K_I$. 所以
\[
  Y=\widehat{K}_I,\qquad e=[K:K_I]=[\widehat{K}:\widehat{K}_I]=[\widehat{K}:Y].
\]
进而, 对每个 $\chi\in Y$, 由 $\chi^*(p)\neq 0$ 可知 $f_\chi\mid n$. 所以 $\chi$ 是由一个确定的模 $n$ 特征 $\chi'$ 所诱导的, 并且 $(\chi')^*=\chi^*$. 另一方面, $\Q(\zeta_n)$ 是 $p$ 在 $\Q(\zeta_m)$ 中的最大不分歧子域, 所以 $K_I\subseteq\Q(\zeta_n)$. 以 $K_D$ 表示 $p$ 在 $K$ 中的分解域, 于是有如下的图表:
\[
  \begin{array}{cc}
    K & \Gal(K/\Q)\\
    \mid & \mid\\
    K_I & \Gal(K_I/\Q)\\
    \mid & \mid\\
    K_D & \Gal(K_D/\Q)\\
    \mid & \mid\\
    \Q & \{1\}
  \end{array}
\]
其中 $I$ 和 $D$ 分别为对应于 $K_I$ 和 $K_D$ 的 $(\Z/n\Z)^\times$ 的子群. 而 $K_I$ 看作 $\Q(\zeta_n)$ 的子域时, 其特征群为 $Y'=\{\chi':\chi\in Y\}$. 由于
\[
  \Gal(K_D/\Q)=(\Z/n\Z)^\times/D\cong\frac{(\Z/n\Z)^\times/I}{D/I},
\]
从而对每个 $\chi'\in Y'=((\Z/n\Z)^\times/I)^{\wedge}$,
\begin{align*}
  \chi'\ \text{为 } \Gal(K_D/\Q)\ \text{的特征}&\Leftrightarrow\chi'(D/I)=1\\
  &\Leftrightarrow\chi'(p)=1\qquad(\text{因为 } D/I\ \text{是由 } \sigma_p\ \text{生成的循环群})\\
  &\Leftrightarrow\chi^*(p)=1\Leftrightarrow\chi\in Z.
\end{align*}
于是 $\widehat{K}_D=Z$, 所以 $f=[K_I:K_D]=[Y:Z]$ 并且 $g=|Z|$. 这就证明了引理.
\end{proof}

现在我们证明: 对每个 Abel 数域 $K$, Dedekind zeta 函数 $\zeta_K(s)$ 是一些对本原特征的 Dirichlet $L$-函数的乘积. 令 $m$ 为域 $K$ 的导子, 则 $K$ 的特征 $\chi$ 都可看成模 $m$ 的 D-特征, 从而 $L(s,\chi^*)$ 是 Dirichlet $L$-函数.

\begin{theorem}[\emph{Hasse}]\label{thm:6.1}
设 $K$ 是 Abel 数域, 则
\[
  \zeta_K(s)=\prod_{\chi\in\widehat{K}}L(s,\chi^*)=\zeta(s)\prod_{\substack{\chi\in\widehat{K}\\ \chi\neq\chi_0}}L(s,\chi^*),
\]
即 Dedekind zeta 函数 $\zeta_K(s)$ 是一些对本原 D-特征的 Dirichlet $L$-函数之乘积.
\end{theorem}

\begin{proof}
根据解析延拓原理, 我们只要对 $\operatorname{Re}(s)>1$ 的情形证明上式即可. 而在 $\operatorname{Re}(s)>1$ 时, 两边均有 Euler 乘积公式: $\zeta_K(s)=\prod_p\prod_{\mathfrak{p}\mid p}(1-N(\mathfrak{p})^{-s})^{-1}$, $L(s,\chi^*)=\prod_p(1-\chi^*(p)p^{-s})^{-1}$. 从而只需对每个有理素数 $p$ 证明
\begin{equation}\tag{1}
  \prod_{\mathfrak{p}\mid p}(1-N(\mathfrak{p})^{-s})=\prod_{\chi\in\widehat{K}}(1-\chi^*(p)p^{-s})
\end{equation}
即可. 设
\[
  pO_K=(\mathfrak{p}_1\cdots\mathfrak{p}_g)^e,\qquad N(\mathfrak{p}_i)=p^f\ (1\leqslant i\leqslant g),\qquad efg=n=[K:\Q],
\]
则 (1) 式左边为 $(1-p^{-fs})^g$. 另一方面, 利用引理 \ref{lem:ch6-1} 和其中的记号, 可知当 $\chi\notin Y$ 时, $\chi^*(p)=0$. 所以 (1) 式右边为 $\prod_{\chi\in Y}(1-\chi^*(p)p^{-s})$. 进而, $\chi\in Z\Leftrightarrow\chi^*(p)=1$. 并且 $Y/Z$ 是 $f$ 阶循环群, 从而 $Y$ 对 $Z$ 的每个陪集中的 $g$ 个特征 $\chi$, $\chi^*(p)$ 取值相同. 并且对不同陪集中的 $f$ 个特征 $\chi_1,\cdots,\chi_f$, 它们对应的本原特征 $\chi_1^*,\cdots,\chi_f^*$ 在 $p$ 的取值分别为 $1,\zeta_f,\cdots,\zeta_f^{f-1}$. 因此
\begin{align*}
  \prod_{\chi\in Y}(1-\chi^*(p)p^{-s})
  &=\prod_{\chi\in Y/Z}(1-\chi^*(p)p^{-s})^g\\
  &=\prod_{i=0}^{f-1}(1-\zeta_f^ip^{-s})^g=(1-p^{-fs})^g.
\end{align*}
这就证明了 (1) 式, 从而也完全证明了定理.
\end{proof}

定理 \ref{thm:4.14} 给出了 $\zeta_K(s)$ 与 $K$ 的理想类数 $h_K$ 的关系. 现在 $\zeta_K(s)$ 可以表成 Dirichlet $L$-函数的乘积, 所以 $h_K$ 可以用 Dirichlet $L$-函数来表达.

\begin{corollarynonum}
设 $K$ 为 Abel 数域, 则
\[
  \rho_Kh_K=\prod_{\substack{\chi\in\widehat{K}\\ \chi\neq\chi_0}}L(1,\chi^*),
\]
其中
\[
  \rho_K=\frac{2^{r_1}(2\pi)^{r_2}R_K}{w_K|d(K)|^{1/2}}.
\]
\end{corollarynonum}

\begin{proof}
定理 \ref{thm:4.14} 给出 $\rho_Kh_K=\lim\limits_{s\to1}\dfrac{\zeta_K(s)}{s-1}$. 而由定理 \ref{thm:6.1} 又知道
\[
  \zeta_K(s)=\prod_{\chi\in\widehat{K}}L(s,\chi^*)=\zeta(s)\prod_{\substack{\chi\in\widehat{K}\\ \chi\neq\chi_0}}L(s,\chi^*),
\]
其中 $\zeta(s)=L(s,\chi_0)$ 是黎曼 zeta 函数. 我们知道当 $\chi\neq\chi_0$ 时, $\chi^*\neq\chi_0$, 从而 $L(s,\chi^*)$ 在 $s=1$ 处正则. 而 $\lim\limits_{s\to1}\dfrac{\zeta(s)}{s-1}=1$. 因此
\[
  \lim_{s\to1}\frac{\zeta_K(s)}{s-1}=\prod_{\substack{\chi\in\widehat{K}\\ \chi\neq\chi_0}}L(1,\chi^*).
\]
这就证明了推论.
\end{proof}

\begin{sloppypar}
现在的问题是将 $L(1,\chi^*)$ 表示成尽可能初等和便于计算的形式. 注意对每个 Dirichlet 特征 $\chi$, $\chi(-1)^2=\chi((-1)^2)=\chi(1)=1$, 从而 $\chi(-1)=\pm 1$. 当 $\chi(-1)=1$ 时, $\chi$ 叫作偶特征, 而 $\chi(-1)=-1$ 时, $\chi$ 叫作奇特征. 今后将看出, 按照 $\chi$ 是奇特征还是偶特征, $L(1,\chi^*)$ 有不同形式的表达式. 首先我们指出特征 $\chi$ 的奇偶性的一个数论意义.
\end{sloppypar}

\begin{lemma}\label{lem:ch6-2}
(a) Abel 数域 $K$ 为实域 $\Leftrightarrow K$ 的特征均为偶特征.

(b) 若 $K$ 为虚 Abel 数域, $K_+$ 为 $K$ 的极大实子域, 则 $[K:K_+]=2$, 并且 $\widehat{K}_+$ 即为 $\widehat{K}$ 中全部偶特征所形成的子群.
\end{lemma}

\begin{proof}
(a) 假设 $\operatorname{cond}(K)=m$, 则有伽罗瓦对应
\[
  \begin{array}{cc}
    \Q(\zeta_m) & (\Z/m\Z)^\times\\
    \mid & \mid\\
    K & G\\
    \mid & \mid\\
    \Q & \{1\}
  \end{array}
\]
于是 $\Gal(K/\Q)=(\Z/m\Z)^\times/G$, 而 $(\Z/m\Z)^\times$ 中的 $-1\ (\mathrm{mod}\ m)$ 相当于复共轭自同构 $\sigma_{-1}(\zeta_m)=\zeta_m^{-1}=\overline{\zeta}_m$. 从而
\begin{align*}
  K\ \text{为实域}&\Leftrightarrow\sigma_{-1}\ \text{为 } K\ \text{中恒等自同构}\Leftrightarrow -1\in G\\
  &\Leftrightarrow -1\ \text{为 } (\Z/m\Z)^\times/G\ \text{中的单位元素}\\
  &\Leftrightarrow K\ \text{的每个特征（即 } (\Z/m\Z)^\times/G\ \text{的每个特征）均为偶特征}.
\end{align*}
(b) 若 $K$ 为虚 Abel 数域, 则 $K$ 中有奇特征, 于是 $\widehat{K}$ 中偶特征全体形成 $\widehat{K}$ 的指数为 $2$ 的子群. 此子群即是 $(\Z/m\Z)^\times/\langle -1,G\rangle$ 的特征群, 而 $\langle -1,G\rangle$ 的固定子域就是 $K$ 的极大实子域 $K_+$, 由此即得结论.
\end{proof}

由引理 \ref{lem:ch6-2} 我们有
\[
  \zeta_{K_+}(s)=\prod_{\chi\in\widehat{K}_+}L(s,\chi^*)=\prod_{\substack{\chi\in\widehat{K}\\ \chi(-1)=1}}L(s,\chi^*)=\zeta(s)\prod_{\substack{\chi\in\widehat{K}\\ \chi(-1)=1\\ \chi\neq\chi_0}}L(s,\chi^*).
\]
从而由定理 \ref{thm:6.1} 的系立刻得到

\begin{lemma}\label{lem:ch6-3}
(a) 设 $K$ 为 $n$ 次实 Abel 域, 则
\[
  R_Kh_K\cdot\frac{2^{n-1}}{|d(K)|^{1/2}}=\prod_{\substack{\chi\in\widehat{K}\\ \chi\neq\chi_0}}L(1,\chi^*).
\]
(b) 设 $K$ 为 $n$ 次虚 Abel 域, 则
\[
  R_Kh_K\frac{(2\pi)^{n/2}}{w_K|d(K)|^{1/2}}=\prod_{\substack{\chi\in\widehat{K}\\ \chi\neq\chi_0}}L(1,\chi^*),
\]
\[
  R_{K_+}h_{K_+}\cdot\frac{2^{n/2-1}}{|d(K_+)|^{1/2}}=\prod_{\substack{\chi_0\neq\chi\in\widehat{K}\\ \chi(-1)=1}}L(1,\chi^*).
\]
\end{lemma}

现在我们来计算 $L(1,\chi^*)$ 的值.

\begin{lemma}\label{lem:ch6-4}
设 $\chi$ 为模 $m$ 本原 D-特征, $m\geqslant 3$, 则
\[
  L(1,\chi)=\begin{cases}
    -\dfrac{2G(1,\chi)}{m}\sum\limits_{1\leqslant k<m/2}\overline{\chi}(k)\log\sin\dfrac{k\pi}{m}, & \text{当 } \chi(-1)=1\ \text{时};\\[1em]
    \dfrac{\pi iG(1,\chi)}{m^2}\sum\limits_{k=1}^{m-1}\overline{\chi}(k)k=\dfrac{\pi iG(1,\chi)}{m(\chi(2)-2)}\sum\limits_{1\leqslant k<m/2}\overline{\chi}(k), & \text{当 } \chi(-1)=-1\ \text{时},
  \end{cases}
\]
其中 $G(k,\chi)=\sum_{a=1}^{m-1}\chi(a)\zeta_m^{ak}$ 为 Gauss 和.
\end{lemma}

\begin{proof}
记 $\omega=\zeta_m$, 当 $\operatorname{Re}(s)>1$ 时
\begin{align*}
  L(s,\chi)&=\sum_{a=0}^{m-1}\chi(a)\sum_{\substack{n=1\\ n\equiv a\ (\mathrm{mod}\ m)}}^{\infty}n^{-s}\\
  &=\sum_{a=0}^{m-1}\chi(a)\sum_{n=1}^{\infty}\left(\frac{1}{m}\sum_{k=0}^{m-1}\omega^{(a-n)k}\right)n^{-s}\\
  &=\frac{1}{m}\sum_{k=0}^{m-1}G(k,\chi)\sum_{n=1}^{\infty}\omega^{-nk}n^{-s}\\
  &=\frac{G(1,\chi)}{m}\sum_{k=0}^{m-1}\overline{\chi}(k)\sum_{n=1}^{\infty}\omega^{-nk}n^{-s}.
\end{align*}
（交换和号是由于级数的绝对收敛性）. 但是 D-级数 $\sum_{n=1}^{\infty}\omega^{-nk}n^{-s}$ 在 $\operatorname{Re}(s)>0$ 中正则, 从而它在 $s=1$ 处连续. 熟知它在 $s=1$ 处的值为 $-\log(1-\omega^{-k})$. 因此
\begin{align*}
  L(1,\chi)&=-\frac{G(1,\chi)}{m}\sum_{k=0}^{m-1}\overline{\chi}(k)\log(1-\omega^{-k})\\
  &=-\frac{\chi(-1)G(1,\chi)}{m}\sum_{k=0}^{m-1}\overline{\chi}(k)\log(1-\omega^k).
\end{align*}
当 $0<k<m$ 时,
\[
  \log(1-\omega^k)=\log 2+\log\sin\frac{k\pi}{m}+\left(\frac{k}{m}-\frac{1}{2}\right)\pi i,
\]
于是（由于 $\operatorname{cond}(\chi)=m\geqslant 3$ 从而 $\chi\neq\chi_0$）
\[
  L(1,\chi)=-\frac{\chi(-1)G(1,\chi)}{m}\sum_{k=1}^{m-1}\overline{\chi}(k)\left(\log\sin\frac{k\pi}{m}+\frac{k\pi i}{m}\right).
\]
当 $\chi(-1)=1$ 时, $\sum\chi(k)k=0$, 而当 $\chi(-1)=-1$ 时,
\[
  \sum\overline{\chi}(k)\log\sin\frac{k\pi}{m}=0
\]
（习题）, 从而
\[
  L(1,\chi)=\begin{cases}
    -\dfrac{2G(1,\chi)}{m}\sum\limits_{1\leqslant k<m/2}\overline{\chi}(k)\log\sin\dfrac{k\pi}{m}, & \chi(-1)=1\ \text{时};\\[1em]
    \dfrac{\pi iG(1,\chi)}{m^2}\sum\limits_{k=1}^{m-1}\overline{\chi}(k)k, & \chi(-1)=-1\ \text{时}.
  \end{cases}
\]
最后再由下面引理 \ref{lem:ch6-5}, 即得到引理 \ref{lem:ch6-4} 的全部结果.
\end{proof}

\begin{lemma}\label{lem:ch6-5}
设 $\chi$ 为模 $m$ 的本原奇特征, 则
\[
  \sum_{k=1}^{m-1}\chi(k)k=\frac{m}{\chi(2)-2}\sum_{1\leqslant k<m/2}\chi(k).
\]
\end{lemma}

\begin{proof}
如果 $2\nmid m$, 则
\begin{align*}
  -\sum_{0<k<m}\chi(k)k&=-\sum_{\substack{0<k<m\\ 2\mid k}}\left(\chi(k)k+\chi(m-k)(m-k)\right)\\
  &=-\sum_{0<k<m/2}\chi(2k)2k+\sum_{0<k<m/2}\chi(2k)(m-2k)\\
  &=\chi(2)\sum_{0<k<m/2}(m-4k)\chi(k).
\end{align*}
因此
\begin{equation}\tag{1}
  -\overline{\chi}(2)\sum_{0<k<m}\chi(k)k=\sum_{0<k<m/2}(m-4k)\chi(k)
\end{equation}
另一方面,
\begin{align*}
  -\sum_{0<k<m}\chi(k)k&=-\sum_{0<k<m/2}\chi(k)k-\sum_{0<k<m/2}\chi(m-k)(m-k)\\
  &=-\sum_{0<k<m/2}\chi(k)k+\sum_{0<k<m/2}\chi(k)(m-k)\\
  &=\sum_{0<k<m/2}(m-2k)\chi(k).
\end{align*}
将此式与 (1) 式合在一起, 即为
\[
  \sum_{0<k<m}\chi(k)k=\frac{m}{\chi(2)-2}\sum_{0<k<m/2}\chi(k).
\]
如果 $4\mid m$, 则 $\chi(2)=0$, 并且 $\chi(k+m/2)=-\chi(k)(0<k<m/2)$（这是因为: 由 $(m/2+1)^2\equiv\frac{m}{4}\cdot m+m+1\equiv 1\ (\mathrm{mod}\ m)$ 可知 $\chi(m/2+1)=\pm 1$. 如果 $\chi(m/2+1)=1$, 则由于对每个奇数 $k$ 均有 $(m/2+1)k\equiv m/2+k\ (\mathrm{mod}\ m)$, 从而 $\chi(m/2+k)=\chi(k)$. 于是导出矛盾 $m=\operatorname{cond}(\chi)\mid m/2$. 因此 $\chi(m/2+1)=-1$, 从而对于每个奇数 $k$ 均有 $\chi(k+m/2)=-\chi(k)$. 而对偶数 $k$ 此式当然也是对的, 因两边均为 $0$）. 于是
\begin{align*}
  \sum_{0<k<m}\chi(k)k&=\sum_{0<k<m/2}\chi(k)k+\sum_{0<k<m/2}\chi(k+m/2)(k+m/2)\\
  &=\sum_{0<k<m/2}\chi(k)k-\sum_{0<k<m/2}\chi(k)(k+m/2)\\
  &=-\frac{m}{2}\sum_{0<k<m/2}\chi(k).
\end{align*}
即引理对 $4\mid m$ 情形也成立.
\end{proof}

将引理 \ref{lem:ch6-4} 代入引理 \ref{lem:ch6-3} 就得到

\begin{lemma}\label{lem:ch6-6}
以 $f_\chi$ 表示 D-特征 $\chi$ 的导子, 则

(a) 对于 $n$ 次实 Abel 数域 $K$,
\[
  R_Kh_K=|d(K)|^{1/2}\prod_{\chi_0\neq\chi\in\widehat{K}}\frac{1}{\sqrt{f_\chi}}\left|\sum_{1\leqslant k<f_\chi/2}\chi(k)\log\sin\frac{k\pi}{f_\chi}\right|.
\]
(b) 对于 $n$ 次虚 Abel 数域 $K$,
\[
  R_{K_+}h_{K_+}=|d(K_+)|^{1/2}\prod_{\substack{\chi_0\neq\chi\in\widehat{K}\\ \chi(-1)=1}}\frac{1}{\sqrt{f_\chi}}\left|\sum_{1\leqslant k<f_\chi/2}\chi(k)\log\sin\frac{k\pi}{f_\chi}\right|,
\]
\[
  \begin{aligned}
    \frac{R_Kh_K}{R_{K_+}h_{K_+}}&=\frac{w_K}{2}|d(K)/d(K_+)|^{1/2}\prod_{\substack{\chi\in\widehat{K}\\ \chi(-1)=-1}}\frac{1}{f_\chi^{3/2}}\left|\sum_{k=1}^{f_\chi-1}\chi(k)k\right|\\
    &=\frac{w_K}{2}|d(K)/d(K_+)|^{1/2}\prod_{\substack{\chi\in\widehat{K}\\ \chi(-1)=-1}}\frac{1}{\sqrt{f_\chi}|\chi(2)-2|}\left|\sum_{1\leqslant k<f_\chi/2}\chi(k)\right|.
  \end{aligned}
\]
\end{lemma}

我们仍然可以将上面诸式右边作进一步化简. 为此, 我们需要第四章所介绍的 $L(s,\chi)$ 和 $\zeta_K(s)$ 的函数方程. 它们可导出 $d(K)$ 和 $f_\chi$ 之间一个简单而奇妙的关系, 就是著名的 Hasse 判别式——导子公式.

\begin{lemma}[\emph{Hasse, 判别式——导子公式}]\label{lem:ch6-7}
设 $K$ 为 Abel 数域, 则
\[
  \prod_{\chi\in\widehat{K}}f_\chi=|d(K)|.
\]
\end{lemma}

\begin{proof}
将公式 $\zeta_K(s)=\prod_{\chi\in\widehat{K}}L(s,\chi^*)$ 两边的 D-级数均应用各自的函数方程. 由 $\zeta_K(s)$ 的函数方程不难算出, 若令
\[
  \Phi_K(s)=\begin{cases}
    |d(K)|^{s/2}(\pi^{-s/2}\Gamma(s/2))^n\zeta_K(s), & \text{若 } K\ \text{为实域};\\[0.5em]
    |d(K)|^{s/2}(\pi^{-s/2}\Gamma(s/2))^{n/2}\left(\pi^{-s/2}\Gamma\left(\frac{1+s}{2}\right)\right)^{n/2}\zeta_K(s), & \text{若 } K\ \text{为虚域},
  \end{cases}
\]
则 $|\Phi_K(s)|=|\Phi_K(1-s)|$. 这里我们使用了公式
\[
  \Gamma(s)=\Gamma(s/2)\Gamma\left(\frac{1+s}{2}\right)\cdot 2^{s-1}/\sqrt{\pi}.
\]
另一方面, 由 $L$-函数的函数方程可知, 若令 $\chi$ 为模 $f_\chi$ 的本原特征, 而
\[
  \Phi(s,\chi)=\begin{cases}
    f_\chi^{s/2}(\pi^{-s/2}\Gamma(s/2))L(s,\chi), & \text{若 } \chi(-1)=1;\\[0.5em]
    f_\chi^{s/2}\left(\pi^{-s/2}\Gamma\left(\frac{s+1}{2}\right)\right)L(s,\chi), & \text{若 } \chi(-1)=-1,
  \end{cases}
\]
则 $|\Phi(s,\chi)|=|\Phi(1-s,\chi)|$. 现在由 $\zeta_K(s)=\prod_{\chi\in\widehat{K}}L(s,\chi^*)$ 即知
\[
  \left|\Phi_K(s)\Big/\prod_{\chi\in\widehat{K}}\Phi(s,\chi)\right|=\left|d(K)\Big/\prod_\chi f_\chi\right|^{s/2}.
\]
但是左边将 $s$ 改为 $1-s$ 时是不变的, 于是右边也应如此. 但这只有 $|d(K)|=\prod f_\chi$ 的时候才可能.
\end{proof}

最后, 将引理 \ref{lem:ch6-7} 和引理 \ref{lem:ch6-6} 放在一起, 我们就给出 Hasse 于 $20$ 世纪 $30$ 年代建立的

\begin{theorem}[\emph{Hasse, Abel 域类数解析公式}]\label{thm:6.2}
(a) 若 $K$ 为 $n$ 次实 Abel 域, 则
\[
  R_Kh_K=\prod_{\chi_0\neq\chi\in\widehat{K}}\left|\sum_{1\leqslant k<f_\chi/2}\chi(k)\log\sin\frac{k\pi}{f_\chi}\right|.
\]
(b) 若 $K$ 为 $n$ 次虚 Abel 域, 则
\[
  R_{K_+}h_{K_+}=\prod_{\substack{\chi_0\neq\chi\in\widehat{K}\\ \chi(-1)=1}}\left|\sum_{1\leqslant k<f_\chi/2}\chi(k)\log\sin\frac{k\pi}{f_\chi}\right|,
\]
\[
  \begin{aligned}
    \frac{R_Kh_K}{R_{K_+}h_{K_+}}&=\frac{w_K}{2}\prod_{\substack{\chi\in\widehat{K}\\ \chi(-1)=-1}}\frac{1}{f_\chi}\left|\sum_{k=1}^{f_\chi-1}\chi(k)k\right|\\
    &=\frac{w_K}{2}\prod_{\substack{\chi\in\widehat{K}\\ \chi(-1)=-1}}\frac{1}{|\chi(2)-2|}\left|\sum_{1\leqslant k<f_\chi/2}\chi(k)\right|.
  \end{aligned}
\]
\end{theorem}

在以下两节里, 我们要对于最早研究类数问题的两种数域——二次域和分圆域给出更简单的类数公式和进一步的结果与猜想.

% ============ 第六章 §2 二次域的类数公式 (原书 pp.217-221) ============
\section{二次域的类数公式}\label{sec:ch6-s2}

根据上节的理论, 为了得到二次域 $K$ 的类数公式, 我们需要决定域 $K$ 的特征群 $\widehat{K}$. 由于 $\Gal(K/\Q)$ 为二元群, $K$ 只有一个非平凡特征, 我们暂时将它记为 $\lambda(\neq\chi_0)$. 显然 $\lambda^2=\chi_0$. 令 $K=\Q(\sqrt{d})$, $d$ 无平方因子. 我们已经证明了 $\operatorname{cond}(K)=|d(K)|$. 因此 $\lambda$ 可看作是模 $|d(K)|$ 的 D-特征. 于是 $\zeta_K(s)=\zeta(s)\cdot L(s,\lambda^*)$. 对于因子 $L(s,\lambda^*)$, 我们需要知道

(a) $\lambda$ 的导子 $f_\lambda$ 是什么？

(b) 能否给出 D-特征 $\lambda$ 的明显表达式？

问题 (a) 容易解决: 利用判别式——导子公式（引理 \ref{lem:ch6-7}）, 对于二次域 $K=\Q(\sqrt{d})$ 则为 $f_\lambda=|d(K)|$. 因此 $\lambda$ 应当是模 $|d(K)|$ 的本原特征. 为了回答问题 (b), 我们需要二次域中的素理想分解定律. 如下定义 $\chi_K(p)$:

(i) 对于奇素数 $p$, 令
\[
  \chi_K(p)=\begin{cases}
    \left(\dfrac{d(K)}{p}\right)=\left(\dfrac{d}{p}\right), & \text{如果 } p\nmid d;\\[0.5em]
    0, & \text{如果 } p\mid d.
  \end{cases}
\]
(ii) $\chi_K(2)=\begin{cases}
  1, & \text{若 } d\equiv 1\ (\mathrm{mod}\ 8);\\
  -1, & \text{若 } d\equiv 5\ (\mathrm{mod}\ 8);\\
  0, & \text{否则}.
\end{cases}$

根据第二章所述, 有理素数 $p$ 在二次域 $K=\Q(\sqrt{d})$ 中分解规律为:

分歧: $p=\mathfrak{p}^2$, 如果 $\chi_K(p)=0$（即 $p\mid d(K)$）;

完全分裂: $p=\mathfrak{p}_1\mathfrak{p}_2$, $\mathfrak{p}_1\neq\mathfrak{p}_2$, 如果 $\chi_K(p)=1$;

惯性: $p=\mathfrak{p}$, 如果 $\chi_K(p)=-1$.

现在我们将 $\chi_K$ 完全积性地将定义域扩充到全体自然数集合上. 换句话说, 如果 $n=p_1^{\alpha_1}\cdots p_s^{\alpha_s}$, 定义 $\chi_K(n)=\chi_K(p_1)^{\alpha_1}\cdots\chi_K(p_s)^{\alpha_s}$. 由于 $|\chi_K(n)|=0$ 或者 $1$, 从而 D-级数 $\sum_{n=1}^{\infty}\chi_K(n)n^{-s}$ 在 $\operatorname{Re}(s)>1$ 中定义出正则函数, 并且有 Euler 展开
\[
  \sum_{n=1}^{\infty}\chi_K(n)n^{-s}=\prod_p(1-\chi_K(p)p^{-s})^{-1}\qquad(\operatorname{Re}(s)>1).
\]

\begin{lemma}\label{lem:ch6-8}
$\lambda=\chi_K$.
\end{lemma}

\begin{proof}
我们已经知道 $\zeta_K(s)=\zeta(s)\sum_{n=1}^{\infty}\lambda(n)n^{-s}$. 根据 D-级数的惟一性定理, 我们只需再证明 $\zeta_K(s)=\zeta(s)\sum_{n=1}^{\infty}\chi_K(n)n^{-s}\ (\operatorname{Re}(s)>1)$ 即可. 由于上式两边均有 Euler 乘积展开, 从而只需对于每个有理素数 $p$ 证明
\begin{equation}\tag{*}
  \prod_{\mathfrak{p}\mid p}(1-N(\mathfrak{p})^{-s})=(1-p^{-s})(1-\chi_K(p)p^{-s})
\end{equation}
即可. 事实上,
\[
  (*)\ \text{式左边}=\begin{cases}
    1-p^{-s}, & (\text{若 } p\ \text{分歧})\\
    (1-p^{-s})^2, & (\text{若 } p\ \text{完全分裂})\\
    (1-p^{-2s}), & (\text{若 } p\ \text{惯性})
  \end{cases}
  =(1-p^{-s})(1-\chi_K(p)p^{-s})=(*)\ \text{式右边}.
\]
于是证明了引理.
\end{proof}

现在很容易给出二次域的类数公式:

\begin{theorem}\label{thm:6.3}
(a) 设 $K$ 为虚二次域并且 $K\neq\Q(\sqrt{-1})$ 和 $\Q(\sqrt{-3})$, 则
\[
  h_K=\frac{1}{|d(K)|}\sum_{k=1}^{|d(K)|}\chi_K(k)k
  =\frac{1}{2-\chi_K(2)}\left|\sum_{1\leqslant k<|d(K)|/2}\chi_K(k)\right|.
\]
(b) 对于实二次域 $K$, 令 $\varepsilon>1$ 为 $K$ 的基本单位, 则
\[
  h_K=\frac{1}{\log\varepsilon}\left|\sum_{1\leqslant k<|d(K)|/2}\chi_K(k)\log\sin\frac{\pi k}{|d(K)|}\right|.
\]
\end{theorem}

\begin{proof}
(a) 对于虚二次域 $K=\Q(\sqrt{-d})$, $d>0$, $r=r_1+r_2-1=0$, 从而 $R_K=1$. 并且当 $|d(K)|>4$ 时, $w_K=2$. （对于 $|d(K)|\leqslant 3$, 则 $K=\Q(\sqrt{-1})$ 和 $\Q(\sqrt{-3})$, 我们在第三章中已经知道它们类数均为 $1$）. 然后由定理 \ref{thm:6.2} 和引理 \ref{lem:ch6-8} 即得 (a) 中类数公式.

(b) 对于实二次域 $K$, $w_K=2$, $r=1$, $R_K=\log\varepsilon$, 从而由定理 \ref{thm:6.2} 和引理 \ref{lem:ch6-8} 即得 (b) 中类数公式.
\end{proof}

\setcounter{example}{0}% 原书例编号分组

\begin{example}
$K=\Q(\sqrt{-5})$, $|d(K)|=20$, $\chi_K(2)=0$, $\chi_K$ 为模 $20$ 的本原 D-特征. 由 $\chi_K$ 的定义算出: $\chi_K$（偶数）$=0$, $\chi_K(1)=\chi_K(9)=1$, $\chi_K(3)=\left(\frac{-5}{3}\right)=1$, $\chi_K(5)=0$, $\chi_K(7)=\left(\frac{-5}{7}\right)=1$. 于是
\[
  h_K=\frac{1}{2}\left|\chi(1)+\chi(3)+\chi(7)+\chi(9)\right|=2.
\]
\end{example}

\begin{example}
$K=\Q(\sqrt{2})$, $|d(K)|=8$, $\chi_K(2)=0$, $\chi_K(1)=1$, $\chi_K(3)=-1$. 基本单位为 $\varepsilon=1+\sqrt{2}$. 于是
\[
  h_K=\frac{1}{\log(1+\sqrt{2})}\left|\chi(1)\log\sin\frac{\pi}{8}+\chi(3)\log\sin\frac{3\pi}{8}\right|
  =\log\left(\frac{\sin\frac{3\pi}{8}}{\sin\frac{\pi}{8}}\right)\Big/\log(1+\sqrt{2}).
\]
即
\[
  (1+\sqrt{2})^{h_K}=\sin\frac{3\pi}{8}\Big/\sin\frac{\pi}{8}=1+2\cos\frac{\pi}{4}=1+\sqrt{2},
\]
从而 $h_K=1$.
\end{example}

在 Боревич 和 Шафаревич 的书《数论》（俄文, $1964$ 年）中列出了 $\Q(\sqrt{-a})$, $1\leqslant a\leqslant 500$ 和 $\Q(\sqrt{d})$, $1\leqslant d\leqslant 500$ 的类数表, 下面是其中的一部分.

（I）虚二次域 $\Q(\sqrt{-a})$, $1\leqslant a\leqslant 100.$

\begin{center}
\begin{tabular}{|c|cccccccccc|}
\hline
$a$ & 1 & 2 & 3 & 5 & 6 & 7 & 10 & 11 & 13 & 14 \\
\hline
$h$ & 1 & 1 & 1 & 2 & 2 & 1 & 2 & 1 & 2 & 4 \\
\hline
$a$ & 15 & 17 & 19 & 21 & 22 & 23 & 26 & 29 & 30 & 31 \\
\hline
$h$ & 2 & 4 & 1 & 4 & 2 & 3 & 6 & 6 & 4 & 3 \\
\hline
\end{tabular}
\end{center}

\begin{center}
\begin{tabular}{|c|cccccccccc|}
\hline
$a$ & 33 & 34 & 35 & 37 & 38 & 39 & 41 & 42 & 43 & 46 \\
\hline
$h$ & 4 & 4 & 2 & 2 & 6 & 4 & 8 & 4 & 1 & 4 \\
\hline
$a$ & 47 & 51 & 53 & 55 & 57 & 58 & 59 & 61 & 62 & 65 \\
\hline
$h$ & 5 & 2 & 6 & 4 & 4 & 2 & 3 & 6 & 8 & 8 \\
\hline
\end{tabular}
\end{center}

\begin{center}
\begin{tabular}{|c|cccccccccc|}
\hline
$a$ & 66 & 67 & 69 & 70 & 71 & 73 & 74 & 77 & 78 & 79 \\
\hline
$h$ & 8 & 1 & 8 & 4 & 7 & 4 & 10 & 8 & 4 & 5 \\
\hline
$a$ & 82 & 83 & 85 & 86 & 87 & 89 & 91 & 93 & 94 & 95 \\
\hline
$h$ & 4 & 3 & 4 & 10 & 6 & 12 & 2 & 4 & 8 & 8 \\
\hline
$a$ & 97 & & & & & & & & & \\
\hline
$h$ & 4 & & & & & & & & & \\
\hline
\end{tabular}
\end{center}

（II）实二次域 $\Q(\sqrt{d})$（$2\leqslant d\leqslant 101$）. $\varepsilon>1$ 表示基本单位, $\omega=\frac{1+\sqrt{d}}{2}$（如果 $d\equiv 1\ (\mathrm{mod}\ 4)$）, $\omega=\sqrt{d}$（如果 $d\equiv 2,3\ (\mathrm{mod}\ 4)$）.

\begin{center}
\begin{tabular}{|c|ccccccccccc|}
\hline
$d$ & 2 & 3 & 5 & 6 & 7 & 10 & 11 & 13 & 14 & 15 & 17 \\
\hline
$h$ & 1 & 1 & 1 & 1 & 1 & 2 & 1 & 1 & 1 & 2 & 1 \\
\hline
$\varepsilon$ & $1+\omega$ & $2+\omega$ & $\omega$ & $5+2\omega$ & $8+3\omega$ & $3+\omega$ & $10+3\omega$ & $1+\omega$ & $15+4\omega$ & $4+\omega$ & $3+2\omega$ \\
\hline
\end{tabular}
\end{center}

\begin{center}
\begin{tabular}{|c|cccccccc|}
\hline
$d$ & 19 & 21 & 22 & 23 & 26 & 29 & 30 & 31 \\
\hline
$h$ & 1 & 1 & 1 & 1 & 1 & 2 & 1 & 2 \\
\hline
$\varepsilon$ & $170+39\omega$ & $2+\omega$ & $197+42\omega$ & $24+5\omega$ & $5+\omega$ & $2+\omega$ & $11+2\omega$ & $1520+273\omega$ \\
\hline
\end{tabular}
\end{center}

\begin{center}
\begin{tabular}{|c|cccccccc|}
\hline
$d$ & 33 & 34 & 35 & 37 & 38 & 39 & 41 & 42 \\
\hline
$h$ & 1 & 2 & 2 & 1 & 1 & 2 & 1 & 1 \\
\hline
$\varepsilon$ & $19+8\omega$ & $35+6\omega$ & $6+\omega$ & $5+2\omega$ & $37+6\omega$ & $25+4\omega$ & $27+10\omega$ & $13+2\omega$ \\
\hline
\end{tabular}
\end{center}

\begin{center}
\begin{tabular}{|c|ccccccc|}
\hline
$d$ & 43 & 46 & 47 & 51 & 53 & 55 & 57 \\
\hline
$h$ & 1 & 1 & 1 & 1 & 2 & 1 & 2 \\
\hline
$\varepsilon$ & $2482+531\omega$ & $24335+3588\omega$ & $48+7\omega$ & $50+7\omega$ & $3+\omega$ & $89+12\omega$ & $131+40\omega$ \\
\hline
\end{tabular}
\end{center}

\begin{center}
\begin{tabular}{|c|ccccccc|}
\hline
$d$ & 58 & 59 & 61 & 62 & 65 & 66 & 67 \\
\hline
$h$ & 2 & 1 & 1 & 1 & 2 & 2 & 1 \\
\hline
$\varepsilon$ & $99+13\omega$ & $530+69\omega$ & $17+5\omega$ & $63+8\omega$ & $7+2\omega$ & $65+8\omega$ & $48842+5967\omega$ \\
\hline
\end{tabular}
\end{center}

\begin{center}
\begin{tabular}{|c|ccccccc|}
\hline
$d$ & 69 & 70 & 71 & 73 & 74 & 77 & 78 \\
\hline
$h$ & 1 & 2 & 1 & 1 & 2 & 1 & 2 \\
\hline
$\varepsilon$ & $11+3\omega$ & $251+30\omega$ & $3480+413\omega$ & $943+250\omega$ & $43+5\omega$ & $4+\omega$ & $53+6\omega$ \\
\hline
\end{tabular}
\end{center}

\begin{center}
\begin{tabular}{|c|ccccccc|}
\hline
$d$ & 79 & 82 & 83 & 85 & 86 & 87 & 89 \\
\hline
$h$ & 3 & 4 & 1 & 2 & 1 & 2 & 1 \\
\hline
$\varepsilon$ & $80+9\omega$ & $9+\omega$ & $82+9\omega$ & $4+\omega$ & $10405+1122\omega$ & $28+3\omega$ & $447+106\omega$ \\
\hline
\end{tabular}
\end{center}

\begin{center}
\begin{tabular}{|c|cccccc|}
\hline
$d$ & 91 & 93 & 94 & 95 & 97 & 101 \\
\hline
$h$ & 2 & 1 & 1 & 2 & 1 & 1 \\
\hline
$\varepsilon$ & $1574+165\omega$ & $13+2\omega$ & $2143295+221064\omega$ & $39+4\omega$ & $5035+1138\omega$ & $9+2\omega$ \\
\hline
\end{tabular}
\end{center}

利用二次域的类数公式, 我们可以对类数的下界给出一些估计（习题 6）. 关于二次域的类数问题, Gauss 有两个著名的猜想:

(1) 只有有限多个虚二次域类数为 $1$;

(2) 存在着无限多个类数为 $1$ 的实二次域.

关于猜想 (1), Gauss 本人计算出当 $K=\Q(\sqrt{-d})$, $d=1,2,3,7,11,19,43,67$ 和 $163$ 时, $h_K=1$. 他预言只有这九个虚二次域的类数为 $1$. $1934$ 年, Heilbronn 证明了猜想 (1), 确切地说, Heilbronn 证明了: 当 $d>163$ 时至多还有一个虚二次域类数为 $1$. 至于这个例外的域是否存在, 是一个长期未解决的问题. 一直到 $1967$ 年才由英国数学家 Baker 和美国数学家 Stark 独立地解决: 那个例外的域是不存在的. 换句话说, 类数为 $1$ 的虚二次域只有 Gauss 发现的那 $9$ 个. 另一方面, 关于 Gauss 猜想 (2), 人们至今未能解决. 实二次域比虚二次域类数问题要困难, 其主要原因是类数公式中多了一个因子 $\log\varepsilon$.

关于二次域类数的估计问题, 对于实二次域
\[
  K=\Q(\sqrt{d})\qquad(d>0),
\]
华罗庚证明了 $h_K<\sqrt{d}$. 这个结果本质上已是最好的可能, 因为在另一方面, 人们证明了: 对于每个 $\varepsilon>0$, 均存在无限多个实二次域 $K=\Q(\sqrt{d})$ 使得 $h_K>d^{\frac{1}{2}-\varepsilon}$. 对于虚二次域 $K=\Q(\sqrt{d})(d<0)$, 类似地证明了当 $|d|>e^{24}$ 时, $h_K\leqslant\frac{1}{3}|d|^{\frac{1}{2}}\log|d|$. 并且对每个 $\varepsilon>0$, 当 $|d|$ 充分大时, 均有 $h_K>|d|^{\frac{1}{2}-\varepsilon}$.

% ============ 第六章 §3 分圆域的类数公式, Kummer 的结果 (原书 pp.222-234) ============
\section{分圆域的类数公式, Kummer 的结果}\label{sec:ch6-s3}

我们已经说过多次, 在历史上是 Kummer 研究 Fermat 猜想导致对分圆域的类数作深入的研究, 这是代数数论的一个重要的源头. Kummer 关于分圆域类数问题的最主要研究结果为

（I）以 $h_p$ 表示分圆域 $\Q(\zeta_p)$ 的类数（$p$ 为奇素数）. 如果 $p\nmid h_p$, 则 Fermat 方程 $x^p+y^p=z^p$ 没有非平凡的有理整数解.

（II）以 $h_p^+$ 表示 $\Q(\zeta_p)$ 的极大实子域 $\Q(\zeta_p+\zeta_p^{-1})$ 的类数, 则 $h_p^+\mid h_p$. 整数 $h_p^+$ 和 $h_p^-=h_p/h_p^+$ 分别叫作是类数 $h_p$ 的第二因子和第一因子.

（III）$p\mid h_p\Leftrightarrow p\mid h_p^-\Leftrightarrow p$ 至少除尽 Bernoulli 数 $B_2,B_4,\cdots,B_{p-3}$ 中一个的分子.

我们在第三章的末尾已经谈到（I）, 并且对于 Fermat 方程的第一种情形给出了证明. 本小节的目的是讲述（II）和（III）, 这就依赖于分圆域的类数公式. 为了简单起见, 我们只考虑分圆域 $K=\Q(\zeta_p)$（$p$ 为奇素数）. 这时 $w_K=2p$. 于是由定理 \ref{thm:6.2} 给出:

\begin{theorem}\label{thm:6.4}
设 $h_p$ 和 $h_p^+$ 分别是 $\Q(\zeta_p)$ 和 $\Q(\zeta_p+\zeta_p^{-1})$ 的类数, $R_p$ 和 $R_p^+$ 分别是它们的 regulator, 则

(a)
\[
  R_p^+h_p^+=\prod_{\substack{\chi\neq\chi_0\\ \chi(-1)=1}}\left|\sum_{k=1}^{(p-1)/2}\chi(k)\log\sin\frac{k\pi}{p}\right|,
\]
其中 $\chi$ 过模 $p$ 的 $\frac{p-3}{2}$ 个非主偶特征.

(b)
\[
  \begin{aligned}
    \frac{R_ph_p}{R_p^+h_p^+}&=p^{-\frac{p-3}{2}}\prod_{\chi(-1)=-1}\left|\sum_{k=1}^{p-1}\chi(k)k\right|\\
    &=p\prod_{\chi(-1)=-1}\frac{1}{|\chi(2)-2|}\left|\sum_{k=1}^{(p-1)/2}\chi(k)\right|,
  \end{aligned}
\]
其中 $\chi$ 过模 $p$ 的 $\frac{p-1}{2}$ 个奇特征.
\end{theorem}

我们在第三章证明了 $R_p=R_p^+\cdot 2^{\frac{p-3}{2}}$, 从而由定理 \ref{thm:6.4} 的 (b) 立刻得到

\begin{corollarynonum}
\[
  h_p^-=(2p)^{-\frac{p-3}{2}}\prod_{\chi(-1)=-1}\left|\sum_{k=1}^{p-1}\chi(k)k\right|
  =2^{-\frac{p-3}{2}}\cdot p\prod_{\chi(-1)=-1}\left(\frac{1}{|\chi(2)-2|}\left|\sum_{k=1}^{(p-1)/2}\chi(k)\right|\right).
\]
\end{corollarynonum}

现在我们用类数公式证明

\begin{theorem}\label{thm:6.5}
$h_p^-\in\Z$.
\end{theorem}

\begin{proof}
模 $p$ 的 $p-1$ 个 D-特征形成一个循环群, 其生成元为 $\chi$, $\chi(g^t)=\zeta_{p-1}^t\ (0\leqslant t\leqslant p-2)$, 其中 $g$ 是模 $p$ 的一个原根. 于是, 模 $p$ 的全部奇特征为 $\{\chi^{2l+1}\mid 0\leqslant l\leqslant\frac{p-3}{2}\}$. 特别地, $\prod_{k=1}^{p-1}\chi(k)k$ 均是分圆域 $\Q(\zeta_{p-1})$ 中的整数, 从而
\[
  B=\prod_{\chi(-1)=-1}\sum_{k=1}^{p-1}\chi(k)k
\]
也是 $\Q(\zeta_{p-1})$ 中整数, 即 $B\in\Z[\zeta_{p-1}]$. 如果我们能够证明 $(2p)^{\frac{p-3}{2}}\mid B$, 则 $(2p)^{-\frac{p-3}{2}}B$ 为整数. 但是
\[
  (2p)^{-\frac{p-3}{2}}B=\pm h_p^-=\pm h_p/h_p^+\in\Q.
\]
这就表明 $h_p/h_p^+\in\Z$, 即 $h_p^-\in\Z$. 因此我们只需证明在 $\Z[\zeta_{p-1}]$ 中 $(2p)^{\frac{p-3}{2}}\mid B$ 即可.

以 $g_s$ 表示 $g^s$ 对于模 $p$ 的最小正剩余, 则
\begin{align*}
  \sum_{k=1}^{p-1}\chi^{2l+1}(k)k&=\sum_{s=0}^{p-2}\chi^{2l+1}(g^s)g_s
  =\sum_{s=0}^{p-2}g_s\zeta_{p-1}^{(2l+1)s}=F(\zeta_{p-1}^{2l+1}),
\end{align*}
其中
\[
  F(x)=\sum_{s=0}^{p-2}g_sx^s\in\Z[x].
\]
于是
\[
  B=F(\zeta_{p-1})F(\zeta_{p-1}^3)\cdots F(\zeta_{p-1}^{p-2}).
\]
我们现在分两步证明 $(2p)^{\frac{p-3}{2}}\mid B$.

（I）证明 $2^{\frac{p-3}{2}}\mid B$: 当 $2\nmid k$ 时, 在 $\Q(\zeta_{p-1})$ 中,
\begin{align*}
  F(\zeta_{p-1}^k)&=\sum_{s=0}^{(p-3)/2}\left(g_s\zeta_{p-1}^{ks}+g_{\frac{p-1}{2}+s}\zeta_{p-1}^{k(\frac{p-1}{2}+s)}\right)\\
  &=\sum_{s=0}^{(p-3)/2}\left(g_s-g_{\frac{p-1}{2}+s}\right)\zeta_{p-1}^{ks}\\
  &=\sum_{s=0}^{(p-3)/2}\zeta_{p-1}^{ks}\left(g_s-(p-g_s)\right)\\
  &\equiv\sum_{s=0}^{(p-3)/2}\zeta_{p-1}^{ks}\pmod{2}.
\end{align*}
从而
\begin{align*}
  F(\zeta_{p-1}^k)(1-\zeta_{p-1}^k)&\equiv\sum_{s=0}^{(p-3)/2}\zeta_{p-1}^{ks}(1-\zeta_{p-1}^k)\\
  &=1-\zeta_{p-1}^{k\cdot\frac{p-1}{2}}=2\equiv 0\pmod{2}.
\end{align*}
于是
\[
  2^{\frac{p-1}{2}}\Big|\prod_{s=0}^{(p-3)/2}F(\zeta_{p-1}^{2s+1})(1-\zeta_{p-1}^{2s+1})
  =B(1-\zeta_{p-1})(1-\zeta_{p-1}^3)\cdots(1-\zeta_{p-1}^{p-2}).
\]
但是
\[
  \prod_{k=1}^{p-2}(1-\zeta_{p-1}^k)=(x^{p-1}-1)'\Big|_{x=1}=p-1;
\]
\[
  \prod_{k=1}^{(p-3)/2}(1-\zeta_{p-1}^{2k})=(x^{\frac{p-1}{2}}-1)'\Big|_{x=1}=\frac{p-1}{2}.
\]
从而 $(1-\zeta_{p-1})(1-\zeta_{p-1}^3)\cdots(1-\zeta_{p-1}^{p-2})=(p-1)\big/\frac{p-1}{2}=2$, 这就证明了 $2^{\frac{p-3}{2}}\mid B$.

（II）证明 $p^{\frac{p-3}{2}}\mid B$: 我们需要考虑 $p$ 在分圆域 $\Q(\zeta_{p-1})$ 中的素理想分解. 由于 $p$ 对于模 $(p-1)$ 的阶为 $1$, 从而由分圆域中素理想分解规律（第二章）可知 $p$ 在 $K=\Q(\zeta_{p-1})$ 中完全分裂, 即
\[
  pO_K=\mathfrak{p}_1\cdots\mathfrak{p}_g,\qquad g=\varphi(p-1).
\]
取 $\mathfrak{p}=\mathfrak{p}_i$, 则

(a) $0,1,\zeta_{p-1},\cdots,\zeta_{p-1}^{p-2}$ 两两模 $\mathfrak{p}$ 不同余（若 $\zeta_{p-1}^i-\zeta_{p-1}^j\in\mathfrak{p}$, $0\leqslant i<j\leqslant p-2$, 则 $1-\zeta_{p-1}^{j-i}\in\mathfrak{p}$. 但是 $\sum_{k=1}^{p-2}(1-\zeta_{p-1}^k)=p-1$, 于是 $p-1\in\mathfrak{p}\mid p$, 这显然不可能）. 由于 $N(\mathfrak{p})=p^f=p$, 从而 $\{0,1,\zeta_{p-1},\zeta_{p-1}^2,\cdots,\zeta_{p-1}^{p-2}\}$ 就是 $\Z[\zeta_{p-1}]$ 模 $\mathfrak{p}$ 的完全代表系.

(b) $1-\zeta_{p-1}^kg\in\mathfrak{p}_i$（对某个 $i$, $1\leqslant i\leqslant g$）$\Leftrightarrow(k,p-1)=1$.

这是因为: 由于 $0\equiv 1-g^{p-1}=\prod_{k=1}^{p-2}(1-\zeta_{p-1}^kg)\ (\mathrm{mod}\ p)$, 从而对于每个 $\mathfrak{p}\mid p$, 均有一个 $1-\zeta_{p-1}^kg$, 使得 $1-\zeta_{p-1}^kg\in\mathfrak{p}$. 根据 (a) 可知, 对于每个固定的 $\mathfrak{p}$, 也只有一个这样的 $1-\zeta_{p-1}^kg$. 进而, 若 $(k,p-1)=d$, 则
\[
  1\equiv\zeta_{p-1}^kg\ (\mathrm{mod}\ \mathfrak{p})\Rightarrow\mathfrak{p}\mid g^{\frac{p-1}{d}}-1\Rightarrow p\mid g^{\frac{p-1}{d}}-1
  \Rightarrow p-1\Big|\frac{p-1}{d}\Rightarrow d=1.
\]
从而若 $1-\zeta_{p-1}^kg\in\mathfrak{p}$, 则必然 $(k,p-1)=1$. 但是满足 $(k,p-1)=1$, $1\leqslant k\leqslant p-1$ 的 $1-\zeta_{p-1}^kg$ 共有 $\varphi(p-1)$ 个, 而 $\mathfrak{p}_i$ 也是 $\varphi(p-1)=g$ 个. 这就表明对于每个 $(k,p-1)=1$ 的 $k$, 均恰好有一个 $p$ 在 $\Q(\zeta_{p-1})$ 中的素理想因子 $\mathfrak{p}$, 使得 $1-\zeta_{p-1}^kg\in\mathfrak{p}$. 我们不妨对 $p$ 的素理想因子重新加以标记, 使得
\[
  (p)=\prod_{\substack{1\leqslant k\leqslant p-1\\ (k,p-1)=1}}\mathfrak{p}_k,\qquad\mathfrak{p}_k\mid(1-\zeta_{p-1}^kg).
\]
现在回到我们的问题上来, 即要证 $p^{\frac{p-3}{2}}\mid B$. 在 $\Z[\zeta_{p-1}]$ 中,
\begin{align*}
  F(\zeta_{p-1}^k)(1-g\zeta_{p-1}^k)&\equiv\sum_{s=0}^{p-2}(g\zeta_{p-1}^k)^s(1-g\zeta_{p-1}^k)\\
  &=1-(g\zeta_{p-1}^k)^{p-1}\equiv 1-g^{p-1}\equiv 0\pmod{p}.
\end{align*}
当 $(k,p-1)>1$ 时, 由 (b) 知对每个 $l$, $(l,p-1)=1$, 均有 $\mathfrak{p}_l\nmid(1-g\zeta_{p-1}^k)$. 从而 $p\mid F(\zeta_{p-1}^k)$. 而当 $(k,p-1)=1$ 时, 如果 $(l,p-1)=1$, $l\neq k$, 则也有 $\mathfrak{p}_l\nmid(1-g\zeta_{p-1}^k)$. 于是
\[
  p\mathfrak{p}_k^{-1}=\prod_{\substack{1\leqslant l\leqslant p-1\\ (l,p-1)=1\\ l\neq k}}\mathfrak{p}_l\Big|F(\zeta_{p-1}^k).
\]
从而 $B=F(\zeta_{p-1})F(\zeta_{p-1}^3)\cdots F(\zeta_{p-1}^{p-2})$ 可被
\[
  p^{\frac{p-1}{2}}\prod_{(k,p-1)=1}\mathfrak{p}_k^{-1}=p^{\frac{p-3}{2}}
\]
除尽. 这就完全证明了定理.
\end{proof}

为了介绍 Kummer 的结果（III）, 我们需要关于 Bernoulli 数的一些重要的性质. 我们在第四章中给出了 Bernoulli 数 $B_n$, 广义 Bernoulli 数 $B_{n,\chi}$ 和 Bernoulli 多项式 $B_n(x)$ 的定义. 它们之间有如下的关系:
\[
  B_n(x)=\sum_{i=0}^{n}\binom{n}{i}B_ix^{n-i},
\]
\[
  B_{n,\chi}=f^{n-1}\sum_{a=1}^{f}\chi(a)B_n\left(\frac{a}{f}\right)\qquad(\chi\ \text{为模 } f\ \text{的 D-特征}),
\]
\[
  B_n=B_{n,\chi_0}\qquad(\text{当 } n\geqslant 2\ \text{时}).
\]

\begin{lemma}\label{lem:ch6-9}
设 $\chi$ 为模 $f$ 的 D-特征.

(a) 对于每个正整数 $F\equiv 0\ (\mathrm{mod}\ f)$, 均有
\[
  B_{n,\chi}=F^{n-1}\sum_{a=1}^{F}\chi(a)B_n\left(\frac{a}{F}\right).
\]
(b) (von Staudt-Clausen) 设 $n$ 为正偶数, 则
\[
  B_n+\sum_{p-1\mid n}\frac{1}{p}\in\Z.
\]
(c) 设 $n$ 为正偶数并且 $n\leqslant p-3$（$p$ 为奇素数）, 则
\[
  \sum_{a=1}^{p-1}a^n\equiv pB_n\pmod{p^2}.
\]
\end{lemma}

\begin{proof}
(a) 根据定义我们有
\[
  \sum_{n=0}^{\infty}B_n(x)\frac{t^n}{n!}=\frac{te^{xt}}{e^t-1},\qquad
  \sum_{n=0}^{\infty}B_{n,\chi}\frac{t^n}{n!}=\sum_{a=1}^{f}\frac{\chi(a)te^{at}}{e^{ft}-1},
\]
从而
\begin{align*}
  \sum_{n=0}^{\infty}F^{n-1}\sum_{a=1}^{F}\chi(a)B_n\left(\frac{a}{F}\right)\frac{t^n}{n!}
  &=\sum_{a=1}^{F}\chi(a)\frac{te^{(a/F)Ft}}{e^{Ft}-1}
  =\sum_{b=1}^{f}\sum_{c=0}^{g-1}\chi(b)\frac{te^{(b+cf)t}}{e^{fgt}-1}\\
  &=\sum_{b=1}^{f}\chi(b)\frac{te^{bt}}{e^{ft}-1}
  =\sum_{n=0}^{\infty}B_{n,\chi}\frac{t^n}{n!},
\end{align*}
由此即得 (a).

(b) 我们只需证明: 对每个素数 $p$ 均有 $p\nmid pB_n$ 的分母, 并且
\begin{equation}\tag{*}
  pB_n\equiv\begin{cases}
    -1\ (\mathrm{mod}\ p), & \text{当 } p-1\mid n\ \text{时};\\
    0\ (\mathrm{mod}\ p), & \text{当 } p-1\nmid n\ \text{时}.
  \end{cases}
\end{equation}
设 $n=2m$. 当 $m=1$ 时, $B_2=\frac{1}{6}=1-\frac{1}{2}-\frac{1}{3}$, 易知 (*) 式对于 $m=1$ 成立. 现在对 $m$ 归纳. 假设 (*) 式对于 $n=2,4,\cdots,2m-2$ 均成立, 由 (a)（$f=1$, $F=p$）, 我们有
\begin{align*}
  pB_{2m}&=pB_{2m,\chi_0}=p^{2m}\sum_{a=1}^{p}B_{2m}\left(\frac{a}{p}\right)\\
  &=p^{2m}\sum_{a=1}^{p}\sum_{j=0}^{2m}\binom{2m}{j}B_j\left(\frac{a}{p}\right)^{2m-j}\\
  &=\sum_{a=1}^{p}\sum_{j=0}^{2m}\binom{2m}{j}(pB_j)a^{2m-j}p^{j-1}\qquad(\text{再利用归纳假设})\\
  &\equiv\sum_{a=1}^{p}\left(pB_0a^{2m}p^{-1}+2mpB_1a^{2m-1}+pB_{2m}p^{2m-1}\right)\ (\mathrm{mod}\ p).
\end{align*}
由于 $B_0=1$, $B_1=-\frac{1}{2}$, $2mpB_1=-mp\equiv 0\ (\mathrm{mod}\ p)$, 于是由上式可知
\[
  (1-p^{2m-1})pB_{2m}\equiv\sum_{a=1}^{p}a^{2m}\equiv
  \begin{cases}
    p-1\pmod{p}, & \text{若 } (p-1)\mid 2m;\\
    0\pmod{p}, & \text{若 } (p-1)\nmid 2m.
  \end{cases}
\]
由于 $1-p^{2m-1}\equiv 1\ (\mathrm{mod}\ p)$, 这就表明 $pB_{2m}\in\Z/p\Z$, 并且
\[
  pB_{2m}\equiv\begin{cases}
    -1\ (\mathrm{mod}\ p), & \text{若 } (p-1)\mid 2m;\\
    0\ (\mathrm{mod}\ p), & \text{若 } (p-1)\nmid 2m.
  \end{cases}
\]

(c) 因为
\begin{align*}
  \sum_{n=0}^{\infty}(B_n(k)-B_n(0))\frac{t^n}{n!}&=\frac{t(e^{kt}-1)}{e^t-1}\\
  &=t(1+e^t+e^{2t}+\cdots+e^{(k-1)t})\\
  &=t+t\sum_{n=0}^{\infty}(1+2^n+\cdots+(k-1)^n)\frac{t^n}{n!}.
\end{align*}
从而当 $n,k\geqslant 1$ 时,
\begin{align*}
  \sum_{a=1}^{k-1}a^n&=\frac{1}{n+1}(B_{n+1}(k)-B_{n+1}(0))\\
  &=\frac{1}{n+1}(B_{n+1}(k)-B_{n+1}).
\end{align*}
从而
\begin{align*}
  \sum_{a=1}^{p-1}a^n&=\frac{1}{n+1}(B_{n+1}(p)-B_{n+1})\\
  &=\frac{1}{n+1}\sum_{k=0}^{n}\binom{n+1}{k}B_kp^{n+1-k}\\
  &\equiv\frac{1}{n+1}\binom{n+1}{n}B_np\equiv pB_n\pmod{p^2}.
\end{align*}
\end{proof}

现在我们讲述 Kummer 的结果（III）.

\begin{theorem}\label{thm:6.6}
$p\mid h_p^-\Leftrightarrow$ 存在 $k\in\{2,4,\cdots,p-3\}$, 使得 $p$ 除尽 $B_k$ 的分子.
\end{theorem}

\begin{proof}
我们在定理 \ref{thm:6.5} 的证明中曾经得到如下的理想恒等式:
\[
  (h_p^-)=\left(\frac{B}{p^{(p-3)/2}}\right)=\prod_{k=1,3,\cdots,p-2}\frac{F(\zeta_{p-1}^k)\mathfrak{p}_k}{p}
\]
（其中 $(k,p-1)>1$ 时规定 $\mathfrak{p}_k=1$）.

于是
\begin{align*}
  p\mid h_p^-=\pm B/p^{(p-3)/2}&\Leftrightarrow\mathfrak{p}_{p-2}\mid p^{-(p-3)/2}B\\
  &\Leftrightarrow\ \text{存在 } k\in\{1,3,\cdots,p-2\},\ \text{使得 } \mathfrak{p}_{p-2}\mid F(\zeta_{p-1}^k)\mathfrak{p}_kp^{-1}\\
  &\qquad(\text{因为 } F(\zeta_{p-1}^k)\mathfrak{p}_k\mathfrak{p}^{-1}\ \text{为整理想})\\
  &\Leftrightarrow\ \text{存在 } k\in\{1,3,\cdots,p-2\},\ \text{使得 } \mathfrak{p}_{p-2}^2\mid F(\zeta_{p-1}^k)\mathfrak{p}_k\\
  &\Leftrightarrow\ \text{存在 } k\in\{1,3,\cdots,p-4\},\ \text{使得 } \mathfrak{p}_{p-2}^2\mid F(\zeta_{p-1}^k)\tag{1}
\end{align*}
最后一式是因为 $\zeta_{p-1}^{p-2}g\equiv 1\ (\mathrm{mod}\ \mathfrak{p}_{p-2})$, 从而
\[
  F(\zeta_{p-1}^{p-2})\equiv\sum_{s=0}^{p-2}(\zeta_{p-1}^{p-2}g)^s\equiv p-1\equiv -1\ (\mathrm{mod}\ \mathfrak{p}_{p-2}),
\]
因此 $\mathfrak{p}_{p-2}^2\nmid F(\zeta_{p-1}^{p-2})\mathfrak{p}_{p-2}$.

现在我们取模 $p$ 原根 $g$, 使得 $g^{p-1}\equiv 1\ (\mathrm{mod}\ p^2)$（设 $g_0$ 为任一模 $p$ 原根, 则 $g_0+lp\ (l\in\Z)$ 均为模 $p$ 原根. 取 $l$ 使得 $g_0^{p-2}l\equiv\frac{g_0^{p-1}-1}{p}\ (\mathrm{mod}\ p)$, 则 $g=g_0+lp$ 即为所求）. 这时
\[
  0\equiv 1-g^{p-1}=\prod_{k=0}^{p-2}(1-\zeta_{p-1}^kg)\ (\mathrm{mod}\ p^2).
\]
与定理 \ref{thm:6.5} 证明中所作的一样, 可以得到: 当 $(k,p-1)=1$ 时, $\mathfrak{p}_k^2\mid 1-\zeta_{p-1}^kg$. 特别取 $k=p-2$, 则 $\zeta_{p-1}\equiv g\ (\mathrm{mod}\ \mathfrak{p}_{p-2}^2)$. 于是
\[
  F(\zeta_{p-1}^k)=\sum_{s=0}^{p-2}g_s\zeta_{p-1}^{sk}\equiv\sum_{s=0}^{p-2}g_sg^{sk}\quad(\mathrm{mod}\ \mathfrak{p}_{p-2}^2).
\]
从而
\begin{equation}\tag{2}
  \mathfrak{p}_{p-2}^2\mid F(\zeta_{p-1}^k)\Leftrightarrow\sum_{s=0}^{p-2}g_sg^{sk}\equiv 0\pmod{p^2}.
\end{equation}
但是 $g_s\equiv g^s+pa_s\ (\mathrm{mod}\ p^2)$, $0\leqslant s\leqslant p-2$, $a_s\in\Z$. 升到 $k+1$ 次幂（$k=1,3,\cdots,p-4$）, 则
\begin{align*}
  g_s^{k+1}&\equiv g^{(k+1)s}+(k+1)g^{sk}pa_s\\
  &\equiv g^{s(k+1)}+(k+1)g^{sk}(g_s-g^s)\\
  &\equiv(k+1)g_sg^{sk}-kg^{s(k+1)}\pmod{p^2}.\tag{3}
\end{align*}
由于 $k+1\leqslant p-3$, 因此 $g^{k+1}\not\equiv 1\ (\mathrm{mod}\ p)$ 而 $g^{p-1}\equiv 1\ (\mathrm{mod}\ p^2)$, 从而
\[
  \sum_{s=0}^{p-2}g^{s(k+1)}=\frac{g^{(p-1)(k+1)}-1}{g^{k+1}-1}\equiv 0\ (\mathrm{mod}\ p^2).
\]
于是由 (3) 式可知
\[
  \sum_{n=0}^{p-1}n^{k+1}=\sum_{s=0}^{p-2}g_s^{k+1}\equiv(k+1)\sum_{s=0}^{p-2}g_sg^{sk}\pmod{p^2}.
\]
再由 (1) 和 (2) 式即知
\[
  p\mid h_p^-\Leftrightarrow\ \text{存在 } k\in\{2,4,\cdots,p-3\},\ \text{使得 } p^2\mid\sum_{n=1}^{p-1}n^k.
\]
但是由引理 \ref{lem:ch6-9} 知道 $p\nmid B_k$ 的分母（因为 $k\leqslant p-3$, $(p-1)\nmid k$）并且 $\sum_{n=1}^{p-1}n^k\equiv pB_k\ (\mathrm{mod}\ p^2)$, 从而 $p^2\mid\sum_{n=1}^{p-1}n^k\Leftrightarrow p\mid B_k$ 的分子. 这就给出最后结果:
\[
  p\mid h_p^-\Leftrightarrow\ \text{存在 } k\in\{2,4,\cdots,p-3\},\ \text{使得 } p\mid B_k\ \text{的分子}.
\]
\end{proof}

\begin{remark}
(1) 定理 \ref{thm:6.6} 相当于 Kummer 结果（III）的一半. 另一半是 $p\mid h_p^-\Leftrightarrow p\mid h_p=h_p^-h_p^+$, 这相当于说: $p\mid h_p^+\Rightarrow p\mid h_p^-$. 为了证明这一点我们需要局部域的进一步的知识. 事实上 Vandiver 猜想: 对每个奇素数 $p$ 均有 $p\nmid h_p^+$. 这个猜想对于 $p<125000$ 均已验证是对的. 但是目前人们还不清楚对于每个奇素数 $p$ 这是否均成立.

(2) 如果 $p\nmid h_p$（即 $p\nmid h_p^-$）, Kummer 称这种 $p$ 为正规素数, 而当 $p\mid h_p$ 时称 $p$ 为不正规素数. Kummer 本人计算了在 $100$ 以内只有 $3$ 个不正规素数: $p=37,59$ 和 $67$. 从 Fermat 猜想的角度看（$p$ 正规则 $x^p+y^p=z^p$ 无非平凡整数解）, 自然希望正规素数愈多愈好. 但是目前人们反而证明了不正规素数有无穷多个, 而正规素数是否有无穷多个, 现在既没有被肯定也没有被否定. 虽然通过大量的计算（Wagstaff 于 $1978$ 年计算出 $125000$ 以内的全部不正规素数）人们倾向于认为正规素数也有无穷多个, 甚至于从概率论上的考虑, 猜想: 正规素数和不正规素数的比例分别各占 $e^{-1/2}\approx 61\%$ 和 $1-e^{-1/2}\approx 39\%$.

(3) Kummer 当年通过手算发现: 当 $p\to+\infty$ 时, $h_p^+$ 增长较慢, 但是 $h_p^-$ 增长飞快. 由于 $h_p=h_p^+h_p^-$, 自然会提出如下的问题:

(A) 类数为 $1$ 的分圆域是否只有有限多个？

(B) 当 $p\to+\infty$ 时, $h_p^+$ 和 $h_p^-$ 的性状如何？

$1976$ 年, Masley 和 Montgomery 完全解决了问题 (A). 他们证明了: 分圆域 $\Q(\zeta_m)$ 的类数 $h_m=1\Leftrightarrow m=3,4,5,7,8,9,11,12,13,15,16,17,19,20,21,24,25,27,28,32,33,35,36,40,44,45,48,60,84$（共 $29$ 个）. 证明中综合性地利用了代数工具和解析工具. Washington 将证明作了更好的整理, 写成书 ``Introduction to Cyclotomic Fields'' 的第 $11$ 章. 差不多与此同时, Masley（$1976$ 年）决定出 $h_m\leqslant 10751$ 的全部 $m(!)$. 如果 $h_m=2\Leftrightarrow m=39,56$; $h_m=3\Leftrightarrow m=23,52,72$ 等等.

至于问题 (B), 关于 $h_p^+$ 的性状是一个很难的数论问题. 其困难主要来自于很难求 $\Q(\zeta_p)$ 的基本单位系, 从而不知道 regulator 值. 目前人们甚至猜不出 $h_p^+$（$p\to+\infty$ 时）的发展规律. Masley 猜想 $h_p^+<p$, 但是 $1982$ 年 Washington 在广义 Riemann 猜想之下证明了对于 $p=11290018777$, 这个猜想是不对的. 对于 $h_p^-$, Kummer 有如下一个著名猜想:
\[
  h_p^-\sim 2p\left(\frac{p}{4\pi^2}\right)^{\frac{p-1}{4}}=2p\left(\frac{p}{39.4784\cdots}\right)^{\frac{p-1}{4}}\qquad(\text{当 } p\to+\infty\ \text{时}).
\]
这一猜想至今未能证明.

下面是对于 $p\leqslant 127$ 的 $h_p^+$ 和 $h_p^-$ 值. （如果广义 Riemann 猜想成立, 可以证明对于每个素数 $p<163$, 均有 $h_p^+=1$. 但是已知 $h_{163}^+=4$.）

\begin{center}
\begin{tabular}{|c|ccccccccccccccc|}
\hline
$p$ & 3 & 5 & 7 & 11 & 13 & 17 & 19 & 23 & 29 & 31 & 37 & 41 & 43 & 47 & 53 \\
\hline
$h_p^+$ & 1 & 1 & 1 & 1 & 1 & 1 & 1 & 1 & 1 & 1 & 1 & 1 & 1 & 1 & 1 \\
\hline
$h_p^-$ & 1 & 1 & 1 & 1 & 1 & 1 & 1 & 3 & 8 & 9 & 37 & $11^2$ & 211 & $5\cdot139$ & 4889 \\
\hline
\end{tabular}
\end{center}

\begin{center}
\begin{tabular}{|c|cccccc|}
\hline
$p$ & 59 & 61 & 67 & 71 & 73 & 79 \\
\hline
$h_p^+$ & 1 & 1 & 1 & 1 & 1 & 1 \\
\hline
$h_p^-$ & $3\cdot59\cdot233$ & $41\cdot1861$ & $67\cdot12739$ & $7^2\cdot79241$ & $89\cdot134353$ & $5\cdot53\cdot377911$ \\
\hline
\end{tabular}
\end{center}

\begin{center}
\begin{tabular}{|c|cccc|}
\hline
$p$ & 83 & 89 & 97 & 101 \\
\hline
$h_p^-$ & $3\cdot279405653$ & $113\cdot118401449$ & $577\cdot3457\cdot206209$ & $5^5\cdot101\cdot601\cdot18701$ \\
\hline
\end{tabular}
\end{center}

\begin{center}
\begin{tabular}{|c|ccc|}
\hline
$p$ & 103 & 107 & 109 \\
\hline
$h_p^-$ & $5\cdot103\cdot1021\cdot17247691$ & $3\cdot743\cdot9859\cdot2886593$ & $17\cdot1009\cdot9431866153$ \\
\hline
\end{tabular}
\end{center}

\begin{center}
\begin{tabular}{|c|cc|}
\hline
$p$ & 113 & 127 \\
\hline
$h_p^-$ & $2^3\cdot17\cdot11853470598257$ & $5\cdot13\cdot43\cdot547\cdot883\cdot3079\cdot626599$ \\
\hline
\end{tabular}
\end{center}

分圆域（及其子域）的类数, 类群和单位群结构, 以及其他数论性质的研究, 一直是代数数论一个富有成果的分支, 并且与模形式理论, 代数几何（特别是椭圆曲线的算术理论）, $p$-adic 分析, 代数 $K$-理论以及群表示理论交织在一起, 构成当前数学研究中极为活跃的边缘性领域. 目前关于分圆域的近代理论已有两种好的参考书（见附录 B）.
\end{remark}

\begin{Exercises}
\item 求证:
\begin{enumerate}[(a)]
  \item 如果 $\chi$ 为模 $m$ 偶特征, 且 $\chi\neq\chi_0$, 则 $\sum_{k=1}^{m-1}\chi(k)k=0$;
  \item 如果 $\chi$ 为模 $m$ 奇特征, 则 $\sum_{k=1}^{m-1}\chi(k)\log\sin\frac{k\pi}{m}=0$.
\end{enumerate}

\item 设 $K$ 是分圆域 $\Q(\zeta_p)$ 的子域（$p$ 为奇素数）, $n=[K:\Q]$, 求证 $|d(K)|=p^{n-1}$. 又若 $K$ 为虚域, 则 $|d(K_+)|=p^{n/2-1}$.

\item (a) 设 $K$ 是分圆域 $\Q(\zeta_p)$ 的二次子域（$p$ 为奇素数）, 则 $K$ 为实二次域 $\Leftrightarrow p\equiv 1\ (\mathrm{mod}\ 4)$;

(b) 更一般地, 设 $K$ 是分圆域 $\Q(\zeta_p)$ 的 $n$ 次子域（$p$ 为奇素数）. 求证 $p\equiv 1\ (\mathrm{mod}\ n)$, 并且: $K$ 为实域 $\Leftrightarrow p\equiv 1\ (\mathrm{mod}\ 2n)$.

\item 用二次域的类数解析公式计算 $\Q(\sqrt{3})$, $\Q(\sqrt{6})$, $\Q(\sqrt{-6})$, $\Q(\sqrt{-23})$ 的类数.

\item 在实二次域 $K$ 中, 求证
\[
  V=\prod_{1\leqslant k<|d(K)|/2}\left(\sin\frac{k\pi}{|d(K)|}\right)^{\chi_K(k)}
\]
是 $K$ 中单位. 进而, 如果以 $U_0$ 表示单位群 $U(K)$ 中由 $V$ 和 $-1$ 生成的乘法子群, 求证: $h(K)=[U(K):U_0]$.

\item 对于虚二次域 $K=\Q(\sqrt{-d})$, $d>3$, $D=|d(K)|$. 求证:
\begin{enumerate}[(a)]
  \item $h(K)\leqslant\frac{1}{D}\sum_{\substack{1\leqslant k<D/2\\ (k,D)=1}}(D-2k)=\frac{\varphi(D)}{2}-\frac{2}{D}\sum_{\substack{1\leqslant k<D/2\\ (k,D)=1}}k$;
  \item 当 $-d\equiv 1\ (\mathrm{mod}\ 4)$ 时, $h(K)<d/4$;

  当 $-d\equiv 2$ 或者 $3\ (\mathrm{mod}\ 4)$ 时, $h(K)<d/2$.
\end{enumerate}

\item 设 $p$ 为奇素数并且 $p\equiv 1\ (\mathrm{mod}\ 4)$. $\varepsilon>1$ 为实二次域 $K=\Q(\sqrt{p})$ 的基本单位. 求证:
\begin{enumerate}[(a)]
  \item $\eta=\prod_{\substack{0<b<p\\ \left(\frac{b}{p}\right)=-1}}\sin\frac{\pi b}{p}\Bigg/\prod_{\substack{0<a<p\\ \left(\frac{a}{p}\right)=1}}\sin\frac{\pi a}{p}$ 为 $K$ 中单位;
  \item $\eta=\varepsilon^{2h(K)}$ 或者 $\eta=\varepsilon^{-2h(K)}$;
  \item $\varepsilon\geqslant\frac{1}{2}(1+\sqrt{5})$;
  \item $\prod_{1\leqslant m<p}\sin\frac{\pi m}{p}=\frac{p}{2^{p-1}}$;
  \item $h(K)<p$.
\end{enumerate}

\item 以 $g_s$ 表示 $g^s$ 模 $p$ 的最小非负剩余, 其中 $p$ 为素数而 $g$ 是模 $p$ 的一个原根, $m=\frac{p-1}{2}$. $h_p^-$ 为分圆域 $\Q(\zeta_p)$ 的类数第一因子. 求证:
\[
  h_p^-=(2p)^{-(m-1)}\left|\det\left(g_{m+i+j}-g_{i+j}\right)_{0\leqslant i,j<m-1}\right|.
\]

\item 计算 $h_7^-$.

\item 证明: $p=37$ 为非正规素数.
\end{Exercises}

